\section{Signal, Data, and Task Definition}

The recognizer receives not an acoustic waveform or video, but a multichannel time series representing tongue--palate contact, and outputs not a closed set of word classes but a character string.

\subsection{EPG Contact Time Series}

EPG records tongue--palate contact during articulation as a multichannel time series and has been used in speech production research, speech training, and articulatory analysis~\cite{hardcastle1989new,hardcastle1990electropalatography,verhoeven2019visualisation}.
In the SmartPalate-style recordings considered here, each time-step frame represents a contact pattern indicating where on the palate contact occurs.
Each trial is treated as a variable-length sequence of EPG contact frames, and the input received by EPGSpeller is this frame sequence.

In the experimental representation used in this paper, SmartPalate contact channels are mapped to a grid aligned with palate layout (a surrogate grid), and aligned contact channels are used.
A surrogate grid assigns each channel to row--column coordinates on the palate, supporting both treating channels as independent scalar features and comparing front ends that exploit spatial layout on the palate.
The physical SmartPalate device specification and the post-export experimental representation used in this paper are not identical.
Table~\ref{tab:signal_summary} specifies the number of aligned channels, the surrogate grid, missing placements, and duplicate placements.

\begin{table}[t]
\centering
\caption{Summary of SmartPalate-style EPG signals (experimental representation).}
\label{tab:signal_summary}
\setlength{\tabcolsep}{2pt}
\renewcommand{\arraystretch}{1.05}
\begin{tabular}{@{}p{0.16\columnwidth}p{0.24\columnwidth}p{0.48\columnwidth}@{}}
\toprule
Item & Value & Role \\
\midrule
Input & Contact frames & Articulatory signal \\
Grid & 16-by-16 surrogate & Spatial front ends \\
Channels & 120 aligned & Multichannel patterns \\
Layout & 4 missing, 5 dup. & Export mismatch \\
Output & Character string & Spelling, not classes \\
Decode & Greedy string & Lexicon-free eval. \\
\bottomrule
\end{tabular}
\end{table}

\subsection{Samples and Labels}

Each sample consists of a variable-length EPG contact sequence and the word label that the participant attempted to spell.
Labels are treated as character strings, not as word classes.
For example, when the target word of a trial is a particular English word, the model learns the character sequence rather than treating the word as a single class identifier.
This design makes EPGSpeller a silent spelling recognizer rather than a finite word classifier.

In a word classifier, the output space is restricted to a vocabulary defined at training or evaluation time.
In a string decoder, the recognizer can output character sequences and can therefore, in principle, compose words that never appeared as a single word class during training.
The evaluation target is how accurately the character string corresponding to the target word can be decoded from the EPG contact time series.

For the primary evaluation, we use character-string output before lexicon projection; lexicon-projected metrics are auxiliary (Table~\ref{tab:task_definition}).

\subsection{Dataset}

This paper uses EPG recordings derived from SilentSpeller~\cite{kimura2022silentspeller}.
The data consist of silent spelling trials recorded with a SmartPalate-style intraoral sensor.
The number of trials, word labels, sequence length, and mean contact rate differ across participants.
These differences are important for interpreting later results.
Within-user recognition measures performance when sufficient calibration data are available; cross-user transfer may manifest differences in contact rate, articulatory style, sensor placement, and oral anatomy as errors.

Table~\ref{tab:dataset_summary} summarizes the dataset for each participant.
p1 and p2 each contain 2328 trials and 1164 word labels.
p3 contains 2797 trials and 1164 word labels, with a lower mean contact rate than the other participants; p4 contains 1167 trials and 1162 word labels with lower repetition density.
Protocol labels P1--P3 denote evaluation splits; participant ids p1--p4 denote SilentSpeller recordings (Table~\ref{tab:dataset_summary}).
Contact rate (Cr.) is the fraction of active palate channels per frame, averaged over trials; lower values indicate sparser contact patterns.
These participant-level differences motivate treating protocols P1/P2 as within-user upper-bound conditions and P3 as the primary transfer bottleneck rather than pooling participants.
Because of p4's lower repetition density, the same calibration conditions as the other participants cannot be set for some within-user evaluations.

\begin{table}[t]
\centering
\caption{Overview of the SilentSpeller-derived dataset (Part.: participant; Tri.: trials; Lab.: word labels; Fr.: mean sequence length in frames; Cr.: mean contact rate).}
\label{tab:dataset_summary}
\setlength{\tabcolsep}{1.5pt}
\renewcommand{\arraystretch}{1.05}
\begin{tabular}{@{}lrrrr@{}}
\toprule
Part. & Tri. & Lab. & Fr. & Cr. \\
\midrule
p1 & 2328 & 1164 & 234 & 0.19 \\
p2 & 2328 & 1164 & 196 & 0.20 \\
p3 & 2797 & 1164 & 283 & 0.12 \\
p4 & 1167 & 1162 & 229 & 0.19 \\
\bottomrule
\end{tabular}
\end{table}
p1/p2 support within-user and calibration evaluation; p3 is the primary transfer and k-shot target; p4 is mainly a transfer target (lower repetition density).

For p1 and p2, repetitions per word are relatively balanced, making it easier to construct within-user unseen-word evaluation and low-shot calibration evaluation.
By contrast, p4 has fewer trials relative to the number of word labels and cannot provide sufficient conditions with multiple repetitions of the same word.
Therefore, p4 is treated mainly as a cross-user transfer target and is not mechanically included in all within-user and calibration conditions.

\subsection{Task Definition and Metrics}
\label{sec:task:definition}

The primary task in this paper is to decode a character string from an EPG contact time series.
The input is a contact frame sequence of varying length.
The output is a normalized alphabetic character string.
The model is not given frame-level character boundaries and learns the correspondence between the contact sequence of an entire trial and the word label.
Accordingly, the recognizer described in later sections uses the connectionist temporal classification (CTC) framework~\cite{graves2006connectionist}.

The primary metric is character error rate (CER) computed from greedy string decoding based on edit distance between the output string and the reference string before lexicon projection.
Lexicon-projected error (LEX) and real-time factor (RTF) are supplementary diagnostics (Table~\ref{tab:task_definition}).

\begin{table}[t]
\centering
\caption{Task definition and primary evaluation.}
\label{tab:task_definition}
\setlength{\tabcolsep}{1.5pt}
\begin{tabular}{@{}p{0.14\columnwidth}p{0.50\columnwidth}p{0.28\columnwidth}@{}}
\toprule
Item & Definition & Role \\
\midrule
Input & EPG contact frames & Not acoustic \\
Target & Character string & Not word class \\
Primary decode & Greedy string & No lexicon \\
Primary metric & CER & Pre-lexicon \\
LEX / RTF & Projection / cost & Supplementary \\
\bottomrule
\end{tabular}
\end{table}

Evaluation conditions are divided into within-user calibration, unseen-word generalization, and cross-user transfer without target calibration; detailed definitions appear in Section~\ref{sec:evaluation} (Table~\ref{tab:protocol_map}).
